\section{Aktivitätsdiagramm}

\begin{tipbox}{Beispiele}
	Das Aktivitätsdiagramm dient der Beschreibung einer \textit{Aktivität}. Das kann z. B. sein
	\begin{itemize}
		\item Ein Algorithmus (z. B. Login-Vorgang in einer Anwendung. Dabei kann der Kontrollfluss und Bedingungen klar dargestellt werden)
		\item Ein Geschäftsprozess (z. B. Bestellprozess in einem Onlineshop. Anwendung findet diese Variante in der Softwareanalyse, denn hier werden Entscheidungen und alternative Abläufe mit verschiedenen Rollen berücksichtigt)
		\item parallele Prozesse (z. B. Bestellbearbeitung in einem Shop-System mit parallelen Zweigen (Zahlung prüfen, Rechnung erzeugen, Lagerbestand aktualisieren, Versand vorbereiten))
	\end{itemize}
\end{tipbox}

\subsection{Übersicht}
\todo{Übersicht Aktivitätsdiagramme}
Beschreibung des Aktivitätsdiagramms...
Das folgende Aktivitätsdiagramm zeigt einen einfachen Prozess, um ein verlorenes Passwort wiederzuerlangen.

\begin{figure}[H]
  \centering
  \includegraphics[width=0.6\textwidth]{\imgpath/act_passwortwiderherstellung.drawio.png}
  \caption{Aktivität Passwortwiederherstellung}
  \label{fig:prozess_passwort_vergessen}
\end{figure}

\subsection{Spezielle Sprachenelemente}
\todo{Sprachenelemente beschreiben ..}

\subsection{Belegsatz}
\begin{figure}[H]
	\centering
	\includegraphics[width=1\textwidth]{\imgpath/ihk_beleg_aktivität.png}
	\caption{IHK-Belegsatz für UML-Aktivitätsdiagramme}
	\label{fig:uml_ihk_belegsatz}
\end{figure} 